% =========================================================================
% SEÇÃO 2: INTRODUÇÃO
% =========================================================================
\section{Introdução}\label{sec:introducao}

A Otimização Combinatória é uma área de pesquisa situada na interseção entre Matemática e Ciência da Computação (veja~\cite{korte2018combinatorial}). Do ponto de vista matemático, a área mantém forte influência com a Teoria dos Grafos, uma vez que digrafos constituem objetos combinatórios com amplo poder de modelagem (veja~\cite{bondy2008graph, diestel2017graph}).
Sob a perspectiva computacional, o foco reside na obtenção de algoritmos eficientes que produzam soluções garantidamente ótimas (veja~\cite{schrijver2003combinatorial, cook1998combinatorial}). Dentre os tópicos de destaque, os problemas de fluxos em digrafos ocupam posição central, tratando de redes onde recursos devem ser transportados respeitando restrições de capacidade (veja~\cite{ahuja1993network}).
O objetivo deste projeto é proporcionar uma formação sólida na área, integrando teoria e prática, abrangendo desde fundamentos conceituais até implementações de alta performance em C++23 e aplicações reais em Visão Computacional.

O problema do $st$-fluxo máximo, onde se busca maximizar a quantidade de fluxo enviada de um vértice fonte $s$ para um vértice de destino $t$, é um dos pilares deste estudo (veja~\cite{ford1956maximal}). Foram analisados e implementados algoritmos clássicos e eficientes, como Ford-Fulkerson, Edmonds-Karp, Dinic e Push-Relabel com Heurística de Gap (veja~\cite{ford1956maximal, edmonds1972theoretical, dinic1970algorithm, goldberg1988new}).
O projeto contempla também o problema de fluxo de custo mínimo, que busca minimizar o custo total enquanto satisfaz as demandas da rede, com destaque para os algoritmos de cancelamento de ciclos de Klein, caminhos sucessivos e o \textit{network simplex}, uma especialização do algoritmo simplex para redes (veja~\cite{dantzig1951application, ahuja1993network}).
A fundamentação teórica desenvolvida reafirma que muitos problemas combinatórios, como o emparelhamento máximo em digrafos bipartidos, coberturas mínimas e segmentação de imagens via corte mínimo, podem ser resolvidos via reduções para problemas de fluxo (veja~\cite{ahuja1993network, boykov2001interactive}).
